\SetPractica{Reacciones de saponificación}{Obtención de un jabón por medio de una reacción de saponificación}

\section{Introducción teórica}



\section{Objetivos}

\begin{itemize}
    \item Obtener un jabón mediante la reacción de una base fuerte, como el hidróxido de potasio, con una grasa animal o un aceite vegetal.
\end{itemize}

\section{Materiales, equipos y reactivos}

\begin{table}[H]
\centering{\small\setstretch{1.25}
\begin{tabular}{|m{0.3\linewidth}|m{0.3\linewidth}|m{0.3\linewidth}|} \hline
    
    \thead{Equipos} & \thead{Materiales} & \thead{Reactivos} \\ \hline
    
    {Balanza \par}
    {Plato caliente}
    
    &
    
    {Vasos de precipitado \par}
    {Matraz Erlenmeyer \par}
    {Agitador \par}
    {Embudo \par}
    {Papel filtro \par}
    {Espátula}
    
    & 
    
    {Hidróxido de potasio (\ce{KOH}) \par}
    {Etanol absoluto (\ce{C2H5OH}) \par}
    {Cloruro de sodio (\ce{NaCl}) \par}
    {Grasa animal o aceite vegetal}
    
    \\ \hline
\end{tabular}}\end{table}

\section{Procedimientos}

\begin{enumerate}
    \item Disolver 9g de hidróxido de potasio en un vaso de precipitado de 100mL una mezcla de 9mL de etanol absoluto y 9mL de agua.
    \item Preparar por separado 40mL de una mezcla etanol-agua al 50\% v/v
    \begin{enumerate}
        \item 20mL de agua y 20mL de etanol
    \end{enumerate}
    
    \item En otro vaso de precipitado de 250mL colocar 5g de grasa animal
    \item Agregar a este último vaso la disolución que contiene hidróxido de potasio.
    \item Durante este tiempo se agrega la disolución de agua-etanol siempre que sea necesario para evitar la formación excesiva de espuma.
    \begin{enumerate}
        \item Calentar la mezcla agitando constantemente, en un baño maría por 30 minutos.
    \end{enumerate}
    
    \item Verter con agitación vigorosa la mezcla anterior en una disolución fría de 25g de cloruro de sodio en 75mL de agua.
    \item Enfriar la mezcla en un baño de hielo por algunos minutos.
    \begin{enumerate}
        \item El jabón o éster metálico se precipita al enfriar.
    \end{enumerate}
    
    \item Filtrar al vacío el precipitado formado, lavándolo 2 veces con porciones de 18mL de agua destilada fría.
    \item Secar al aire por 24 horas, pesar el precipitado y calcular el rendimiento.
\end{enumerate}